\begin{definition}
    Пусть $T\neq\emptyset$. Тогда $\R^T$~-- множество всех вещественных функций на $T$, то есть $\{x : T \to \R\}$
\end{definition}

Тогда понятно, что $X_{\bullet}(\omega) \in \R^T$

\begin{definition}
    \textit{Пространство траекторий}: $\R^T$
\end{definition}

\begin{definition}
    \textit{Элементарный цилиндр}: $C_{t, A} = \{x\in\R^T : x(t)\in A\}$, где $A$~- борелевсое на $\R$.
\end{definition}

\begin{note}
    То есть множество траекторий (или функций), которые в момент времени $t$ проходят через  $A$.
\end{note}

\begin{example}
    Если $T=\{1,2\}$, $t=1$, $A=[a,b]$, то $C_t$~-- цилиндр (полоса) с основанием $[a, b]$ на первой оси.
\end{example}

\begin{statement}
     Элементарные цилиндры порождают алгебру цилиндров вида: 
     \begin{equation*}
         C_{t_1, \ldots, t_n, A} = \{x\in\R^T : \Big(x(t_1), \ldots, x(t_n)\Big)\in A\} \text{, \; где $A$~-- борелевсое на $\R^n$.}
     \end{equation*}
\end{statement}

\begin{proof}
    Чтобы увидеть, что в самом деле это алгебра, достаточно заметить, что два цилиндра с множествами точек $t_1, \ldots , t_n$ и $s_1, \ldots, s_k$ и основаниями $A \subset \R^n$ и $B \subset \R^k$ могут быть записаны с помощью объединенного множества точек $(t_1, \ldots , t_n, s_1, \ldots , s_k)$ и оснований в общем пространстве $\R^{n+k}$, полученных умножением исходных оснований на $\R^k$ и $\R^n$. соответственно.
\end{proof}

\begin{note}
    Представление цилиндра $C_{t_1, \ldots, t_n, A}$ не однозначно, всегда можно добавить лишнюю точку $t_{n+1}$ и заменить $A$ на $A\times\R^1$, т.е. реальной зависимости от значений в $t_{n+1}$ не будет,  но можно выбрать наименьшее возможное число точек $t_j$ в представлении. Для цилиндра, отличного от всего $\R^T$, точки в таком минимальном наборе определены однозначно, дальнейшая неединственность представления связана только с их возможной перенумерацией.
\end{note}

\begin{note}
    Цилиндры еще не составляют $\sigma$-алгебру
\end{note}

\begin{definition}
    \textit{Цилиндрическая $\sigma$-алгебра} $\mathcal{B}_T$~-- это наименьшая $\sigma$-алгебра, содержащая все цилиндры (это то же самое, что наименьшая $\sigma$-алгебра, содержащая все элементарные цилиндры).
\end{definition}

\begin{note}
    В терминах $\sigma$-алгебр, порожденных функциями, можно сказать, что $\mathcal{B}_T$ есть $\sigma$-алгебра в пространстве траекторий $\R^T$, порожденная всеми функциями вида $x \mapsto x(t), \; t \in T$
\end{note}